% \iffalse meta-comment
%
% hit-float.dtx 是浮动体与图表
%
% \fi
%
% \section{浮动体与图表}
%
% ^^A ======================================================================
% ^^A Module thesis-floats: 浮动体默认行为与双语题注（hit-thesis）
% ^^A ======================================================================
%    \begin{macrocode}
%<*thesis-floats>
%    \end{macrocode}
%
% \changes{v3.0.0}{2020/05/21}{添加博后的浮动间距设置}
% \changes{v3.2a}{0000/00/00}{浮动体间距的三个分支迁到 expl3}
% 博后单独一套间距；其余按 \option{style/glue} 决定要不要给伸缩量。
%    \begin{macrocode}
\bool_if:NTF \g__hit_postdoc_bool
  {
    \skip_set:Nn \intextsep    { 10.5bp }
    \skip_set:Nn \textfloatsep { 10.5bp }
    \skip_set:Nn \floatsep     { 10.5bp }
  }
  {
    \bool_if:NTF \g__hit_glue_bool
      {
        \skip_set:Nn \intextsep    { 8.50398bp~plus~2.83465bp~minus~0bp }
        \skip_set:Nn \textfloatsep { 8.50398bp~plus~2.83465bp~minus~0bp }
        \skip_set:Nn \floatsep     { 12bp~plus~2.83465bp~minus~0bp }
      }
      {
        \skip_set:Nn \intextsep    { 8.50398bp }
        \skip_set:Nn \textfloatsep { 8.50398bp }
        \skip_set:Nn \floatsep     { 12bp }
      }
  }
%    \end{macrocode}
%
% 此处设置float在p选项时间隔，此处不设置\cs{@fptop}和\cs{@fpbot}以确保居中。
% \changes{v1.0.12}{2018/04/03}{修正float为p状态时默认不居中bug}
% \changes{v2.0.04}{2018/12/04}{删除\cs{@fpsep}设置，似乎没有什么用}
% \changes{v2.0.04}{2018/12/04}{更新\cs{intextsep}\cs{textfloatsep}\cs{floatsep}间距为正文行间距}
% 下面这组命令使浮动对象的缺省值稍微宽松一点，从而防止幅度对象占据过多的文本页面，
% 也可以防止在很大空白的浮动页上放置很小的图形。
% \changes{v1.0.08}{2017/11/05}{修改附录中图、表、公式数字编码}
%    \begin{macrocode}
\tl_gput_right:Nn \appendix { \cs_set_nopar:Npn \thefigure { \thechapter - \arabic {figure} } }
\tl_gput_right:Nn \appendix { \cs_set_nopar:Npn \thetable { \thechapter - \arabic {table} } }
\tl_gput_right:Nn \appendix { \cs_set_nopar:Npn \theequation { \thechapter - \arabic {equation} } }
\cs_set:Npn \textfraction      { 0.15 }
\cs_set:Npn \topfraction       { 0.85 }
\cs_set:Npn \bottomfraction    { 0.65 }
\cs_set:Npn \floatpagefraction { 0.60 }
%    \end{macrocode}
%
% 由于我工的双标题，导致标题之下多出一空白字符的距离，去除。
% \changes{v2.0.04}{2018/12/04}{更新图段后空白距离}
% \changes{v2.0.04}{2018/12/04}{删除表段后空白距离}
% \changes{v2.0.05}{2018/12/05}{删除图段后空白距离}
%    \begin{macro}{\@makecaption}
% 根据我工规范，本科和硕博的图题序号之后的空格不一样。
%    \begin{hitrgu}[\PGR][2.13.1]
% 每个图均应有图题（由图序和图名组成），图题不宜有标点符号，图名在图序之后空1个
% 半角字符排写。
%    \end{hitrgu}
%    \begin{hitrgu}[\UGR][2.13.1]
% 每个图均应有图题（由图序和图名组成），图题不宜有标点符号，图名在图序之后空1个
% 字符排写。
%    \end{hitrgu}
% 我工规范中没有明确规定是否标题是否居中对齐，这里给出一个居中选项自行调整。
% 注意，我工只规定：“居中书写”。此处不额外添加悬挂处理。
% \changes{v1.0.06}{2017/10/25}{此处更改了选项的名称}
% \changes{v1.0.07}{2017/11/04}{优化了最后一行居中算法，使其两边对齐、单词内部断行}
%
% \changes{v3.2a}{0000/00/00}{\cs{@makecaption} 与 \cs{LT@makecaption} 迁到 expl3}
% \cs{hit@caption@separator} 是题号与题名之间的间隔，本科一个汉字宽，其余是 \cs{enskip}。
% \cs{hit@caption@last@line} 把最后一行单独取出来居中，\option{capcenterlast} 开时用。
% 里面的盒子编号与分组照抄原来的写法，只是换成 expl3 的拼法。量宽仍用 \cs{sbox}
% 而不是 \cs{hbox\_set:Nn}：前者会套一层 \cs{color@begingroup}，盒子里的颜色
% 改动不会漏出去，后者没有。
%    \begin{macrocode}
\cs_new:Npn \hit@caption@separator { \hit@if@option@TF { bachelor } { \hskip \ccwd } { \enskip } }
\cs_new_protected:Npn \hit@caption@last@line #1
  {
    \vskip 6.3bp
    {
      \tex_setbox:D 0 = \tex_vbox:D {#1}
      \tex_setbox:D 1 = \tex_vbox:D
        {
          \tex_unvbox:D 0
          \tex_setbox:D 2 = \tex_lastbox:D
          \hbox_to_wd:nn { \textwidth }
            { \hfill \tex_unhcopy:D 2 \unskip \unskip \hfill }
        }
      \tex_unvbox:D 1
    }
  }
\cs_set:Npn \@makecaption #1#2
  {
    \vskip \abovecaptionskip
    \wuhao
    \sbox \@tempboxa { #1 \hit@caption@separator #2 }
    \dim_compare:nNnTF { \box_wd:N \@tempboxa } > { \hsize }
      {
        \hit@if@option@TF { capcenterlast }
          { \hit@caption@last@line { #1 \hit@caption@separator #2 } }
          { #1 \hit@caption@separator #2 }
        \par
      }
      {
        \tex_global:D \@minipagefalse
        \hb@xt@ \hsize { \hfil \box \@tempboxa \hfil }
      }
    \vskip \belowcaptionskip
  }
%    \end{macrocode}
%
%    \end{macro}
%    \begin{macro}{\LT@makecaption}
% 有关longtable的caption支持
%    \begin{macrocode}
\cs_set:Npn \LT@makecaption #1#2#3
  {
    \LT@mcol \LT@cols c
      {
        \hbox_to_zero:n
          {
            \hss
            \parbox [t] \LTcapwidth
              {
                \wuhao
                \sbox \@tempboxa { #1 { #2 \hit@caption@separator } #3 }
                \dim_compare:nNnTF { \box_wd:N \@tempboxa } > { \hsize }
                  {
                    \hit@if@option@TF { capcenterlast }
                      { \hit@caption@last@line { #1 { #2 \hit@caption@separator } #3 } }
                      { #1 { #2 \hit@caption@separator } #3 }
                    \par
                  }
                  {
                    \tex_global:D \@minipagefalse
                    \hbox_to_wd:nn { \hsize } { \hfil \box \@tempboxa \hfil }
                  }
                \endgraf \vskip \belowcaptionskip
              }
            \hss
          }
      }
  }
%    \end{macrocode}
%
%    \end{macro}
%    \begin{macro}{\longbionenumcaption}
% 长表格的双语标题是一个坑. 因为第一不能用浮动格式，只能用longtable包中的tabular
% ，这样表题只能使用表格中前两行来写。这样出现了一个问题是，中英表题的间距，标题
% 和表第一行间距，表格内部间距等多个变量的协调问题。这个问题只要使用tabular的形
% 式，就是无解的。唯一的方法就是把这些参数都给用户列出来。以下，第2，5参数为中英
% 双语标题内容，1，4为标题参数。6为中英标题间距，7为表题和表格间距。
%    \begin{macrocode}
\cs_set_nopar:Npn \longbionenumcaption #1 #2 #3 #4 #5 #6 #7
  {
%
\@if@contemptyarg{#1}{\caption{#2}}{\caption[#1]{#2}}%
\cs_gset_eq:NN \@cont@oldtablename \tablename
\cs_gset_nopar:Npn \tablename {#3}~
\cs_gset_eq:NN \LT@c@ption \@cont@LT@nonumintoc
\\[#6]~
\@if@contemptyarg{#4}{\caption{#5}}{\caption[#4]{#5}}%
\cs_gset_eq:NN \tablename \@cont@oldtablename
\cs_gset_eq:NN \LT@c@ption \@cont@oldLT@c@ption
\vspace{#7}
  }
%    \end{macrocode}
%
%    \end{macro}
%    \begin{macro}{\ltfontsize}
% \changes{v3.2a}{0000/00/00}{长表格字号机制抽到 \file{src/utils/hit-table.dtx}}
% 我们采用 \pkg{longtable} 来处理跨页的表格。默认字号压到五号，机制见
% \file{src/utils/hit-table.dtx}。
%    \begin{macrocode}
\hit@setup@table
%    \end{macrocode}
%
%    \end{macro}
% 图表名称及格式。
% \changes{v1.0.08}{2017/11/05}{删除冗余公式符号定义}
%    \begin{macrocode}
\cs_set:Npn \thesubtable { ( \alph { subtable } ) }
\cs_set:Npn \thefigure   { \arabic { chapter } - \arabic { figure } }% 使图编号为 7-1 的格式
%    \end{macrocode}
%
% \changes{v3.0.15}{2020/05/21}{根据窝工2021“四月革命”之精神，修改子图为双括号}
%    \begin{macrocode}
\cs_set:Npn \thesubfigure { ( \alph { subfigure } ) }% 子图编号为 a) 的格式
\cs_set:Npn \p@subfigure  { \thefigure \nobreakspace { } }% 子图引用为 7-1 a)
\cs_set:Npn \thetable     { \arabic { chapter } - \arabic { table } }% 表编号为 7-1
%    \end{macrocode}
%
% \changes{v3.2a}{0000/00/00}{\cs{setitemize}/\cs{setenumerate} 从 book-floats 迁出至 book-paragraph}
% 此处删除hang caption的设置
% \changes{v1.0.06}{2017/10/25}{删除caption hang 的默认设置，因为不在规范要求中}
%    \begin{macrocode}
\captionnamefont{\wuhao}
\captiontitlefont{\wuhao}
\cs_set:Npn \subcapsize { \wuhao }
%    \end{macrocode}
%
% \option{modern} 下子图题的字体、标签分隔与最后一行对齐交给 \pkg{caption}，
% 键名不同但取的是同一批值。
%    \begin{macrocode}
\DeclareCaptionFont { hit@subcap } { \subcapsize }
\DeclareCaptionLabelSeparator { hit@sub } { \hskip \subfiglabelskip }
\captionsetup [ sub ]
  { font = hit@subcap , labelformat = simple , labelsep = hit@sub , skip = 0pt }
\bool_if:NT \g__hit_subcapcenterlast_bool
  { \captionsetup [ sub ] { justification = centerlast } }
\dim_set:Nn \abovecaptionskip {0pt}%为了双标题之间的间距，不能设置
\dim_set:Nn \belowcaptionskip {0pt}
%    \end{macrocode}
%
% \changes{v3.2a}{0000/00/00}{\env{publist} 挪到 \file{src/utils/hit-publist.dtx}，各个类共用}
% \env{publist} 见 \file{src/utils/hit-publist.dtx}。
%    \begin{macrocode}
\hit@setup@publist
%    \end{macrocode}
%
% \changes{v2.0.01}{2018/06/28}{去除定理注释括号}
% \changes{v3.2a}{0000/00/00}{ntheorem 定理样式与 \cs{arraycolsep} 从 book-floats 迁出至 book-math；\cs{setlength}\{\cs{parindent}\}\{2em\} 迁出至 book-paragraph}
%
% \changes{v3.2a}{0000/00/00}{原 hithesis.sty 中的旋转/算法/代码段宏包并入 book-floats}
% 旋转图表所需的 \pkg{rotating}。
%    \begin{macrocode}
\RequirePackage{rotating}
%    \end{macrocode}
%
% 算法宏包。注意加载顺序为 float、hyperref、algorithm2e，否则无法生成算法列表。
% 算法标题居中、中英双语标题、按章节编号等设置参见原文档。
%    \begin{macrocode}
\RequirePackage[algoruled,linesnumbered,algochapter]{algorithm2e}
\SetAlCapSty{}
\NewDocumentCommand \foocaption { m }
  {~\cs_set:Npn \@algocf@pre@plainruled
    {
      \hrule height 1.5pt depth 0pt \kern \interspacetitleruled #1~\kern \interspacealgoruled
      \hrule height 1pt depth 0pt \kern \interspacetitleruled
    }~}
\cs_set:Npn \@algocf@post@ruled
  { \kern \interspacealgoruled \hrule height 1.5pt \relax }
\cs_new:Npn \algoenname { Algo. }
\newfloatlist[chapter]{algoen}{aen}{\listalgoenname}{\algoenname}
\newfixedcaption{\algoencaption}{algoen}
\cs_set:Npn \thealgoen { \thechapter - \arabic { algocf } }
\cs_set:Npn \@cftmakeaentitle
  { \chapter * { \listalgoenname \@mkboth { \listalgoenname } { \listalgoenname } } }
\cs_set:Npn \algorithmcfname { 算法 }
\dim_set:Nn \AlCapSkip {1.2ex}
\SetAlgoSkip{1pt}
\RenewDocumentCommand \algocf@captiontext { m m }
  { \wuhao #1 \algocf@typo \nobreakspace { } \AlCapFnt { } #2 }
\cs_if_exist:NTF \algocf@within
  {
    \cs_set:Npn \thealgocf
      { \use:c { the \algocf@within } - \@arabic \c@algocf }
  }
  { \cs_set:Npn \thealgocf { \@arabic \c@algocf } }
\RenewDocumentCommand \algocf@makecaption { m m }
  {
    \dim_add:Nn \hsize { \algomargin }
    \sbox \@tempboxa { \algocf@captiontext {#1} {#2} }
    \hskip .5\algomargin
    \parbox [t] { \hsize } { \centering \algocf@captiontext {#1} {#2} }
    \dim_add:Nn \hsize { -\algomargin }
  }
\NewDocumentCommand \AlgoBiCaption { m m }
  {
    \caption [#1]
      { \protect \dim_set:Nn \baselineskip { 1.5em } #1~\protect~\\ Algo.\thealgocf \nobreakspace { } #2 }
    \addcontentsline { aen } { algoen } { \protect \numberline { \thealgoen } {#2} }
  }
%    \end{macrocode}
%
% 代码段排版 \pkg{listings}（默认柔性列宽与自动断行）。
%    \begin{macrocode}
\RequirePackage{listings}
\lstset{
columns=flexible,
breaklines=true
}
%    \end{macrocode}
%
% \changes{v3.2a}{0000/00/00}{原 hithesis.sty 中的示例 tikz 样式与染色体示意图绘制宏并入 book-floats}
% 示例文档使用的 tikz 样式与染色体段落、细胞绘制宏。
%    \begin{macrocode}
\tikzset{maternal/.style={colorone}}
\tikzset{paternal/.style={colortwo}}
\tikzset{variant/.style={colorthree!80!colorzero}}
\tikzset{reference/.style={colorzero}}
\tikzset{aallele/.style={colorzero,rotate=90}}
\tikzset{ballele/.style={colorthree!80!colorzero,rotate=90}}
\tikzset{refseg/.style={colorzero,draw=colorzero,~opacity=0.2}}
\tikzset{mseg/.style={colorone,draw=colorone,~opacity=0.2}}
\tikzset{pseg/.style={colortwo,draw=colortwo,~opacity=0.2}}
\tikzset{vseg/.style={colorthree!80!colorzero,draw=colorthree!80!colorzero,~opacity=0.6}}
\tikzset{bncell/.style={draw=colorzero,opacity=0.2,line~width=2pt,~rounded~corners=1pt}}
\tikzset{btcell/.style={draw=colorone,opacity=0.6,~line~width=2pt,~rounded~corners=1pt}}
\tikzset{tncell/.style={colorzero,opacity=0.9}}
\tikzset{ttcell/.style={colorone,opacity=0.6}}
\tikzset{tscell/.style={colorzero}}
\tikzset{refcell/.style={colorzero}}
\tikzset{evolve/.style={->,draw=colortwo,opacity=0.3,line~width=1.5pt}}
\tikzset{fakeevolve/.style={->,draw=colorzero,opacity=0.3,line~width=1.5pt}}
\tikzset{refline/.style={dashed,draw=colorzero,line~width=1pt}}
\tikzset{tnline/.style={dashed,draw=colorzero,opacity=0.3,line~width=1pt}}
\NewDocumentCommand \gseg { m m m m m m m m m }
  {
%
~\pgfmathsetmacro{\sstartx}{#1}~
 \pgfmathsetmacro{\slengx}{#2}~
 \pgfmathsetmacro{\sy}{#3}~
 \pgfmathsetmacro{\sdy}{#4}~
 \pgfmathsetmacro{\sdx}{#5}~
 \pgfmathsetmacro{\sdxh}{#7}~
 \pgfmathsetmacro{\sdxt}{#8}~
 \fill[#6](\sstartx,\sy)--(\sstartx-\sdx,\sy+\sdy)--~
 (\slengx+\sstartx+1.5-\sdx,\sy+\sdy)--(\slengx+\sstartx+1.5,\sy)--~
 (\slengx+\sstartx+1.5-\sdx,\sy-\sdy)--(\sstartx-\sdx,\sy-\sdy)--cycle;~
 \draw[#9]~(\sstartx-\sdxh,\sy)~--~(\sstartx,~\sy);~
 \draw[#9]~(\slengx+\sstartx+1.5,~\sy)~--~(\slengx+\sstartx+1.5+\sdxt,\sy);~
  }
\NewDocumentCommand \gsegr { m m m m m m m m m }
  {
%
~\pgfmathsetmacro{\sstartx}{#1}~
 \pgfmathsetmacro{\slengx}{#2}~
 \pgfmathsetmacro{\sy}{#3}~
 \pgfmathsetmacro{\sdy}{#4}~
 \pgfmathsetmacro{\sdx}{#5}~
 \pgfmathsetmacro{\sdxh}{#7}~
 \pgfmathsetmacro{\sdxt}{#8}~
 \fill[#6](\sstartx-0.5,\sy)--(\sstartx+\sdx-0.5,\sy+\sdy)--~
 (\slengx+\sstartx+1.5+\sdx-0.5,\sy+\sdy)--(\slengx+\sstartx+1.5-0.5,\sy)--~
 (\slengx+\sstartx+1.5+\sdx-0.5,\sy-\sdy)--(\sstartx+\sdx-0.5,\sy-\sdy)--cycle;~
 \draw[#9]~(\sstartx-\sdxh-0.5,\sy)~--~(\sstartx-0.5,~\sy);~
 \draw[#9]~(\slengx+\sstartx+1.5-0.5,~\sy)~--~(\slengx+\sstartx+1.5+\sdxt-0.5,\sy);~
  }
\NewDocumentCommand \rcell { m m }
  {
%
~\pgfmathsetmacro{\x}{#1}~
 \pgfmathsetmacro{\y}{#2}~
 \draw (\x+1,\y)~node[aallele]{A};~
 \draw (\x+2,\y)~node[aallele]{C};~
 \draw (\x+3,\y)~node[aallele]{T};~
 \draw (\x+4,\y)~node[aallele]{C};~
 \gseg{\x}{4}{\y}{0.2}{0.5}{refseg}{1.5}{1}{reference}~
  }
\NewDocumentCommand \ncell { m m }
  {
%
~\pgfmathsetmacro{\x}{#1}~
 \pgfmathsetmacro{\y}{#2}~
 \draw[bncell](\x-2,\y+0.5)--(\x+7,\y+0.5)--~
 (\x+7,\y-1)--(\x-2,\y-1)--cycle;~
 \draw (\x+1,\y)~node[aallele]{A};~
 \draw (\x+2,\y)~node[ballele]{G};~
 \draw (\x+3,\y)~node[aallele]{T};~
 \draw (\x+4,\y)~node[aallele]{C};~
 \gseg{\x}{4}{\y}{0.2}{0.5}{mseg}{1.5}{1}{maternal}~
 \draw (\x+1,\y-0.5)~node[ballele]{T};~
 \draw (\x+2,\y-0.5)~node[aallele]{C};~
 \draw (\x+3,\y-0.5)~node[aallele]{T};~
 \draw (\x+4,\y-0.5)~node[ballele]{A};~
 \gseg{\x}{4}{\y-0.5}{0.2}{0.5}{pseg}{1.5}{1}{paternal}~
  }
\NewDocumentCommand \tcellone { m m }
  {
%
~\pgfmathsetmacro{\x}{#1}~
 \pgfmathsetmacro{\y}{#2}~
 \draw[btcell](\x-2,\y+0.5)--(\x+7,\y+0.5)--~
 (\x+7,\y-1.5)--(\x-2,\y-1.5)--cycle;~
 \draw (\x+1,\y)~node[aallele]{A};~
 \draw (\x+2,\y)~node[ballele]{G};~
 \draw (\x+3,\y)~node[aallele]{T};~
 \draw (\x+4,\y)~node[aallele]{C};~
 \gseg{\x}{4}{\y}{0.2}{0.5}{mseg}{1.5}{1}{maternal}~
 \draw (\x+1,\y-0.5)~node[aallele]{A};~
 \draw (\x+2,\y-0.5)~node[ballele]{G};~
 \draw (\x+3,\y-0.5)~node[aallele]{T};~
 \draw (\x+4,\y-0.5)~node[aallele]{C};~
 \gseg{\x}{4}{\y-0.5}{0.2}{0.5}{mseg}{1.5}{1}{maternal}~
 \draw (\x+1,\y-1)~node[ballele]{T};~
 \draw (\x+2,\y-1)~node[aallele]{C};~
 \draw (\x+3,\y-1)~node[aallele]{T};~
 \draw (\x+4,\y-1)~node[ballele]{A};~
 \gseg{\x}{4}{\y-1}{0.2}{0.5}{pseg}{1.5}{1}{paternal}~
  }
\NewDocumentCommand \tcellthree { m m }
  {
%
~\pgfmathsetmacro{\x}{#1}~
 \pgfmathsetmacro{\y}{#2}~
 \draw[btcell](\x-2,\y+0.5)--(\x+11,\y+0.5)--~
 (\x+11,\y-1)--(\x-2,\y-1)--cycle;~
 \draw (\x+1,\y)~node[aallele]{A};~
 \draw (\x+2,\y)~node[ballele]{G};~
 \gseg{\x}{2}{\y}{0.2}{0.5}{mseg}{1.5}{0}{maternal}~
 \gseg{\x+4}{0}{\y}{0.2}{0.5}{vseg}{0.5}{0.5}{variant}~
 \draw (\x+7,\y)~node[aallele]{T};~
 \draw (\x+8,\y)~node[aallele]{C};~
 \gseg{\x+6}{2}{\y}{0.2}{0.5}{mseg}{0}{1}{maternal}~
 \draw (\x+1,\y-0.5)~node[ballele]{T};~
 \draw (\x+2,\y-0.5)~node[aallele]{C};~
 \draw (\x+3,\y-0.5)~node[aallele]{T};~
 \draw (\x+4,\y-0.5)~node[ballele]{A};~
 \gseg{\x}{4}{\y-0.5}{0.2}{0.5}{pseg}{1.5}{1}{paternal}~
  }
\NewDocumentCommand \tcellfour { m m }
  {
%
~\pgfmathsetmacro{\x}{#1}~
 \pgfmathsetmacro{\y}{#2}~
 \draw[btcell](\x-2,\y+0.5)--(\x+15,\y+0.5)--~
 (\x+15,\y-1)--(\x-2,\y-1)--cycle;~
 \draw (\x+1,\y)~node[aallele]{A};~
 \draw (\x+2,\y)~node[ballele]{G};~
 \gseg{\x}{2}{\y}{0.2}{0.5}{mseg}{1.5}{0}{maternal}~
 \gseg{\x+4}{0}{\y}{0.2}{0.5}{vseg}{0.5}{0.5}{variant}~
 \draw (\x+7,\y)~node[aallele]{T};~
 \gseg{\x+6}{1}{\y}{0.2}{0.5}{mseg}{0}{0}{maternal}~
 \gseg{\x+9}{0}{\y}{0.2}{0.5}{vseg}{0.5}{0.5}{variant}~
 \draw (\x+12,\y)~node[aallele]{C};~
 \gseg{\x+11}{1}{\y}{0.2}{0.5}{mseg}{0}{1}{maternal}~
 \draw (\x+1,\y-0.5)~node[ballele]{T};~
 \draw (\x+2,\y-0.5)~node[aallele]{C};~
 \draw (\x+3,\y-0.5)~node[aallele]{T};~
 \draw (\x+4,\y-0.5)~node[ballele]{A};~
 \gseg{\x}{4}{\y-0.5}{0.2}{0.5}{pseg}{1.5}{1}{paternal}~
  }
\NewDocumentCommand \tcelltwo { m m }
  {
%
~\pgfmathsetmacro{\x}{#1}~
 \pgfmathsetmacro{\y}{#2}~
 \draw[btcell](\x-2,\y+0.5)--(\x+7,\y+0.5)--~
 (\x+7,\y-1.5)--(\x-2,\y-1.5)--cycle;~
 \draw (\x+1,\y)~node[aallele]{A};~
 \draw (\x+2,\y)~node[ballele]{G};~
 \draw (\x+3,\y)~node[aallele]{T};~
 \draw (\x+4,\y)~node[aallele]{C};~
 \gseg{\x}{4}{\y}{0.2}{0.5}{mseg}{1.5}{1}{maternal}~
 \draw (\x+1,\y-0.5)~node[aallele]{A};~
 \draw (\x+2,\y-0.5)~node[ballele]{G};~
 \draw (\x+3,\y-0.5)~node[aallele]{T};~
 \draw (\x+4,\y-0.5)~node[ballele]{G};~
 \gseg{\x}{4}{\y-0.5}{0.2}{0.5}{mseg}{1.5}{1}{maternal}~
 \draw (\x+1,\y-1)~node[ballele]{T};~
 \draw (\x+2,\y-1)~node[aallele]{C};~
 \draw (\x+3,\y-1)~node[aallele]{T};~
 \draw (\x+4,\y-1)~node[ballele]{A};~
 \gseg{\x}{4}{\y-1}{0.2}{0.5}{pseg}{1.5}{1}{paternal}~
  }
\NewDocumentCommand \tcellfive { m m }
  {
%
~\pgfmathsetmacro{\x}{#1}~
 \pgfmathsetmacro{\y}{#2}~
 \draw[btcell](\x-2,\y+0.5)--(\x+9.5,\y+0.5)--~
 (\x+9.5,\y-1.5)--(\x-2,\y-1.5)--cycle;~
 \draw (\x+1,\y)~node[aallele]{A};~
 \draw (\x+2,\y)~node[ballele]{G};~
 \draw (\x+3,\y)~node[aallele]{T};~
 \draw (\x+4,\y)~node[aallele]{C};~
 \gseg{\x}{4}{\y}{0.2}{0.5}{mseg}{1.5}{1}{maternal}~
 \draw (\x+1,\y-0.5)~node[aallele]{A};~
 \draw (\x+2,\y-0.5)~node[ballele]{G};~
 \draw (\x+3,\y-0.5)~node[aallele]{T};~
 \draw (\x+4,\y-0.5)~node[aallele]{C};~
 \gseg{\x}{4}{\y-0.5}{0.2}{0.5}{mseg}{1.5}{1}{maternal}~
 \draw (\x+1,\y-1)~node[ballele]{T};~
 \gseg{\x}{1}{\y-1}{0.2}{0.5}{pseg}{1.5}{0}{paternal}~
 \draw (\x+4.5,\y-1)~node[ballele]{A};~
 \draw (\x+5.5,\y-1)~node[aallele]{T};~
 \draw (\x+6.5,\y-1)~node[aallele]{C};~
 \gsegr{\x+3.5}{3}{\y-1}{0.2}{0.5}{pseg}{0.5}{1.5}{paternal}~
  }
%</thesis-floats>
%    \end{macrocode}
%
% \Finale
